%% Detailed OptEvolve loop (TikZ picture only).
%% Compile through flowmap-standalone.tex to the PDF included by the paper.
%% Coordinate map, in cm: input y=12.5; genome x=0.2..7.4;
%% gate/compile/measure/route x=8.54..15.76 (a 1.14-cm arrow channel).
%% The colored stages on the right have 0.45-cm vertical gaps; the input
%% has the same gap to the blue genome and amber gate. The blue bands share
%% x=0.52..7.08 (0.32-cm side margins) and have 0.42-cm vertical gaps.
%% Gene pool top y=2.89 gives ~0.5 cm clearance below both columns.
%% This is Figure 1's palette and visual grammar, expanded to eight layers.
\begingroup%
\definecolor{oeInput}{HTML}{E7E0F1}% supplied problem
\definecolor{oeGenome}{HTML}{DCE6F7}% typed solver design and compilation
\definecolor{oeGate}{HTML}{F4DFA5}% pre-execution check
\definecolor{oeOutput}{HTML}{E3EFD7}% measurement and routing
\definecolor{oePool}{HTML}{E5E9EA}% retained evidence
\definecolor{oeInk}{HTML}{293846}% outlines and flow
\definecolor{oePrice}{HTML}{B3242B}% measured-cost feedback only
\begin{tikzpicture}[
 x=1cm,y=1cm,font=\sffamily,line cap=round,line join=round,
 panel/.style={draw=oeInk,line width=.8pt,rounded corners=2.3mm,
               align=center,inner sep=1.5mm},
 band/.style={draw=oeInk!65,line width=.6pt,rounded corners=1.2mm,
              fill=white,inner sep=1mm},
 heading/.style={font=\sffamily\fontsize{9.5}{11}\bfseries\selectfont},
 group/.style={font=\sffamily\fontsize{8.5}{10}\bfseries\selectfont,
               anchor=west},
 row/.style={font=\sffamily\fontsize{8.2}{9.7}\selectfont,anchor=west},
 copy/.style={font=\sffamily\fontsize{8.2}{10}\selectfont,align=center},
 flow/.style={-{Latex[length=2.3mm]},draw=oeInk,line width=.95pt},
 evidence/.style={-{Latex[length=2.3mm]},draw=oeInk!72,line width=.85pt,
                   dash pattern=on 2mm off 1.1mm},
 price/.style={-{Latex[length=2.4mm]},draw=oePrice,line width=1pt}
]
%% Inputs banner: the family is the unit searched; the device informs the price.
\node[panel,fill=oeInput,minimum width=16.1cm,minimum height=1.0cm]
  (input) at (8.20,12.50) {};
\node[heading] at (8.20,12.72) {Problem family + structure};
\node[copy] at (8.20,12.29)
  {instances at different scales \quad $\cdot$ \quad target accuracy \quad $\cdot$ \quad CPU / GPU profile};

%% Genome container: three bands share the same width and eight numbered rows.
\node[panel,fill=oeGenome,minimum width=7.20cm,minimum height=8.15cm]
  (genome) at (3.80,7.47) {};
\node[heading] at (3.80,11.10) {Typed genome: change one layer at a time};

%% Form band: choose the algorithm family and how a problem family is traversed.
%% All bands use the same width and left text inset, so their edges line up.
\node[band,minimum width=6.56cm,minimum height=1.36cm]
  at (3.80,9.92) {};
\node[group] at (.78,10.26) {FORM};
\node[row] at (.78,9.90) {1\; Family: splitting / path / working set};
\node[row] at (.78,9.53) {2\; Continuation: grid / interpolation};

%% Algorithm band: state reuse, safe restriction, handoff and numerics.
\node[band,minimum width=6.56cm,minimum height=2.60cm]
  at (3.80,7.52) {};
\node[group] at (.78,8.58) {ALGORITHM};
\node[row] at (.78,8.21) {3\; State: warm start / factorization update};
\node[row] at (.78,7.77) {4\; Restrict: screen + re-admit};
\node[row] (handoff) at (.78,7.33) {5\; Handoff: when to switch algorithms};
\node[row] at (.78,6.89) {6\; Numerics: precision / fast math};

%% Hardware band: admission and parallel width change per-step cost.
\node[band,minimum width=6.56cm,minimum height=1.42cm]
  at (3.80,5.09) {};
\node[group] at (.78,5.46) {HARDWARE};
\node[row] at (.78,5.09) {7\; Admission: work versus startup cost};
\node[row] at (.78,4.70) {8\; Width: GPU tile / CPU threads};

%% Stage-0 gate: family-specific obligations; the proof in this paper
%% formally covers the splitting schemes and relaxation rules.
\node[panel,fill=oeGate,minimum width=7.22cm,minimum height=2.10cm]
  (gate) at (12.15,10.49) {};
\node[heading] at (12.15,11.12) {Stage-0: check before compiling};
\node[copy,text width=6.55cm] at (12.15,10.28)
  {Splitting: convergence conditions\\
   Exact path: face + step \quad $\cdot$ \quad Working set: full KKT pass};

%% Execution remains blue: compilation/fusion instantiates the genome's
%% device gene after admission; it is not another correctness gate.
\node[panel,fill=oeGenome,minimum width=7.22cm,minimum height=1.55cm]
  (compile) at (12.15,8.215) {};
\node[heading] at (12.15,8.515) {Compile + fuse};
\node[copy] at (12.15,7.965) {choose a GPU tile or CPU thread width};

%% Evaluation is green like the routed output: both belong to Figure 1's
%% Measure and route stage. The red line explains the return to layer 5.
\node[panel,fill=oeOutput,minimum width=7.22cm,minimum height=1.75cm]
  (measure) at (12.15,6.115) {};
\node[heading] at (12.15,6.48) {Measure to acceptance};
\node[copy] at (12.15,6.05)
  {$T=\mathrm{setup}+N\,t$ across shapes and devices};
\node[copy,text=oePrice] at (12.15,5.60)
  {measured $t$ re-prices the layer-5 handoff};

%% Result: the fastest accepted route can be an existing specialist.
\node[panel,fill=oeOutput,minimum width=7.22cm,minimum height=1.45cm]
  (route) at (12.15,4.065) {};
\node[heading] at (12.15,4.355) {Solver + routing rule};
\node[copy] at (12.15,3.805) {evolved or specialist; checked answer};

%% Evidence pool: success and failure both affect later proposals. The
%% dashed line goes around the left margin to avoid crossing the genome.
\node[panel,fill=oePool,minimum width=16.1cm,minimum height=1.08cm]
  (pool) at (8.20,2.35) {};
\node[heading] at (8.20,2.57) {Gene pool: measured moves, positive and negative};
\node[copy] at (8.20,2.10)
  {Keep the evidence and its code dependency; invalidate it when the code changes.};

%% Main path: input -> genome -> gate -> compilation -> measurement -> route.
\draw[flow] (3.80,12.0) -- (genome.north);
\draw[flow] (genome.east |- gate.center) -- (gate.west);
\draw[flow] (gate.south) -- (compile.north);
\draw[flow] (compile.south) -- (measure.north);
\draw[flow] (measure.south) -- (route.north);

%% Dashed return: measurements update the gene pool, which proposes one
%% typed change to the next genome. Keep this separate from the red price arc.
\draw[evidence] (measure.east) -- (16.55,6.115) --
  (16.55,2.35) -- (pool.east);
\draw[evidence] (pool.west) -- (-.32,2.35) --
  (-.32,6.14) -- (.20,6.14);

%% Red return: the per-step cost t(op,device), measured at the right,
%% moves the crossover in the layer-5 handoff. Terminate on the actual
%% handoff row, not the outer panel, so the dependency is unambiguous.
\draw[price] (measure.west) -- (7.78,6.115) --
  (7.78,7.33) -- (handoff.east);
\end{tikzpicture}%
\endgroup%
